\section{What a dark node believes: baseline, systematic, residual}
\label{sec:grbelieve}

Everything so far solves \emph{one instant}: whoever is lit right now
speaks, and the universe rearranges.  But a trading day is not an
instant, and the sharpest test of the machinery is a node that was
seen --- briefly, tellingly --- and has been dark since.  This section
works that test on the smallest example that contains it, and the
answer will organize everything a dark node knows into three parts.

\subsection{A morning, replayed by hand}

Two assets.  $A$ is liquid and prints every hour: its handle reads
$10,11,12,13,14,15$ at hours $0$ through $5$ (toy numbers, in handle
points).  $B$ is thin: it prints $10$ at hour $0$ and then only once
more, $10$ again, at hour $3.5$.  The desk believes $B$ follows $A$
one for one --- amplitude one --- and believes the influence runs one
way: $A$ is the market, $B$ a satellite.  What should $B$'s published
mark do?

Walk it as a trader would.  From hour $0$ to $3$, $B$ is dark and
nothing distinguishes it from its $A$-driven prediction, so it rides
the relation: $10,11,12,13$.  At hour $3.5$, $B$ prints $10$ while the
freshest $A$ information says $13$.  The print is a fact; publish it.
But it also \emph{teaches} something: relative to what $A$ implies,
$B$ is dislocated by $10-13=-3$ points.  When $A$ ticks to $14$ at
hour $4$ and $B$ is dark again, the desk does not throw that lesson
away: the mark should be $14-3=11$ --- following $A$'s move at full
amplitude \emph{from the dislocated level}.  The blue staircase of
\cref{fig:grstory}(a) is this path.

Now hand the same tape to the machinery of
\cref{sec:grmeet} run afresh at each instant, with whoever is printing
right at that moment as the lit set.  (The story is walked in raw
levels for readability; against a flat baseline of $10$, a level is
its innovation plus ten, and with amplitude one the two solves differ
by that constant shift and nothing else.)  The dashed violet path is
what comes back, and it is wrong twice, in instructive ways.  At hour
$3.5$, only $B$ is lit --- and the solve drags $A$ down to
\MacGrStorySnapA: the contract read backwards is exactly as precise as
read forwards (\cref{sec:gredge}), and with nothing else lit, one
satellite print legitimately marks down the index.  At hour $4$, only
$A$ is lit --- and $B$ snaps to \MacGrStorySnapB: the solve holds no
record that $B$ ever disagreed, so the dislocation is erased by the
next source tick, where the desk path stands at \MacGrStoryDeskB.
Nothing is broken in the algebra.  The solver answered exactly the
question it was asked --- ``what do the contracts and \emph{this
instant's} prints imply?'' --- and the desk was asking a different
question, one whose answer needs two things a per-instant solve does
not have: \emph{direction} and \emph{memory}.

\begin{figure}[htbp]
\centering
\paperfig{\textwidth}{fig_gr_story.pdf}
\caption{The asynchronous pair.  (a)~$A$ prints hourly (ink); $B$
prints twice (orange squares).  The desk path (blue) rides the
relation between prints, learns the $-3$-point dislocation at the
$3.5$-hour print, and keeps it: $14-3=11$ at hour $4$.  The
per-instant solve (dashed violet) drags $A$ to the print at hour
$3.5$ (no direction) and erases the dislocation at hour $4$ (no
memory).  (b)~The carried dislocation under three half-lives: what
$t_{1/2}$ asserts is how long a private dislocation stays believed.}
\label{fig:grstory}
\end{figure}

\subsection{No amplitude can fake a memory}

Before adding machinery, close one tempting escape.  After the print,
could a softer \emph{amplitude} give the desk a gentler path than the
staircase?  Following $A$ at a constant fraction $\beta$ from the
printed level publishes, at hour $t$,
\[
B=10+\beta\,\bigl(A_t-13\bigr):
\]
a fixed share of $A$'s move since the print, forever.  Now suppose
the desk's instinct is that the share should \emph{change} as the
print goes stale --- wishing, say, for half of $A$'s move by hour $4$
but three quarters of it by hour $5$: marks of $10.5$ and $11.5$.
Match each wish separately:
\[
10+\beta\,(14-13)=10.5
\;\Longrightarrow\;
\beta=0.5,
\qquad
10+\beta\,(15-13)=11.5
\;\Longrightarrow\;
\beta=0.75 .
\]
No constant amplitude grants both: a share that varies with elapsed
time is not an amplitude at all --- it is a \emph{state} of the node,
evolving on its own clock.  The coherent version of ``the print goes
stale'' is a dislocation that decays: publish
$B=A_t-3\cdot2^{-\text{elapsed}/t_{1/2}}$, the remembered $-3$
fading with a stated half-life.  Such paths start on the staircase
and drift \emph{up} toward the full relation as the memory fades ---
and that direction is forced.  At hour $4$, for every half-life,
\[
14-3\cdot2^{-0.5/t_{1/2}}\;\ge\;14-3=11 :
\]
forgetting a dislocation can only close the gap toward $A$, never
deepen it.  So no half-life --- and, by the contradiction above, no
amplitude either --- holds the hour-$4$ mark below $11$; the
half-way wish of $10.5$ is thereby diagnosed rather than granted,
since it asks to undo part of a move the desk itself asserted at
amplitude one.  The menu after a print is exactly one dial wide:
from the flat staircase (never forget) to an instant snap-back
(forget at once), parametrized by the half-life.

\subsection{The residual state}

So give each node the state the story demands.  Write $\mathcal{U}_i$
for node $i$'s \emph{persisted dislocation}: the part of its last
actual observation that its informers did not explain, measured at the
print as the gap between the print and the contemporaneous
systematic prediction ($10-13=-3$ in the story).  Between
observations the dislocation is believed with a stated memory,
\begin{equation}
\mathcal{U}_i(\text{now})
=\mathcal{U}_i(\text{print})\cdot
2^{-\Delta t_{\mathrm{days}}/t_{1/2}},
\label{eq:grdecay}
\end{equation}
a half-life decay on the elapsed time, with elapsed time and
half-life quoted on one clock --- the universe construction uses
days, the story above ran in hours.  (And the \emph{variance} of the
carried value grows as it ages: the longer since the print, the less
certain what remains of it.)  Panel~(b) of \cref{fig:grstory} draws
the decay for three settings of the dial on the story's hourly clock:
half an hour (a dislocation is short-lived noise), two hours, and
never-forget; the staged universe of \cref{sec:grcomplete} runs the
same dial in days, its sister name carrying a two-day half-life.  The
half-life is not estimated from anything in this book --- it is a
stated desk commitment, exactly like the anchor $\kappa$, and the
audit of \cref{sec:grcomplete} is what holds it honest.

How does the state enter the solve?  Through the contract, as an
offset: a node carrying a dislocation obeys its relations \emph{from
the dislocated level}, so the factor becomes
\[
p_{ij}\,\bigl(\Theta_i-\beta_{ij}\Theta_j-\mathcal{U}_i\bigr)^{2},
\]
with the carried $\mathcal{U}_i$ from \cref{eq:grdecay} and its grown
variance widening the relation noise.  Everything else --- the
assembly, the accounting, the anchors --- is unchanged.

Direction is the remaining assertion, and it is worth stating exactly
because the algebra cannot supply it.  When $B$'s surprising print
arrived, the desk revised \emph{B's} state only --- it did not mark
the index down.  That is a structural claim about the world
(influence flows one way), and \cref{sec:gredge} showed a symmetric
contract cannot express it: tight trust binds both directions.  The
reference implementation therefore cuts the update at the source.
Concretely: a satellite's print enters the machinery as a measurement
of its \emph{own} residual --- the gap between the print and the
contemporaneous systematic prediction is booked into
$\mathcal{U}$ --- and not as an observation row of \cref{eq:grjoint},
which the joint solve would propagate backward into the informer.
This is a stated policy layered on top of the Gaussian, not a
consequence of it.  This book keeps to the policy's one load-bearing sentence: a
dark universe must never let its thinnest member mark down its
benchmark.

\subsection{Three priors, one belief}

Step back and collect what a single dark node now knows, because it
has become a complete inventory.  Its published handle is
\[
\text{baseline}
\;+\;
\underbrace{\text{systematic}}_{\text{the graph's message}}
\;+\;
\underbrace{\text{residual}}_{\text{its own memory}},
\]
three parts of decreasing evidence, each with its own uncertainty and
its own clock.  The \emph{baseline} is Chapter~10's: yesterday's
posterior, transported by Chapter~9, its variance grown on
Chapter~8's clock --- what the node would publish if the universe said
nothing.  The \emph{systematic} part is this chapter's: what the lit
universe's innovations imply through the contracts, at the amplitudes
asserted, shrunk by the anchor, priced by the accountant.  The
\emph{residual} is the node's own remembered dislocation, decaying on
its half-life.  Nothing else enters, and each part is separately
auditable --- which is the property that lets \cref{sec:grcomplete}
print, for any published mark, exactly where every point of it came
from.  The machinery is complete; it remains to run it.
